\documentclass{beamer}
\usepackage[utf8]{inputenc}
\usepackage[spanish]{babel}
\usepackage{amsmath}
\usepackage{graphicx}
\usepackage{listings}

\usetheme{Madrid}
\usecolortheme{whale}

\title[Pipelines y Buenas Prácticas]{Módulo 4: Aprendizaje Automático}
\subtitle{Clase 6: Mini Proyecto y Arquitectura de Pipelines}
\author{Diplomado en Ciencia de Datos}
\date{\today}

\begin{document}

\begin{frame}
    \titlepage
\end{frame}

% ---------------------------------------------------
\begin{frame}{Objetivos de la Sesión}
    \begin{itemize}
        \item Evolucionar de la programación algorítmica a la \textbf{Ingeniería de Modelos (MLOps)}.
        \item Entender el peligro del \textit{Data Leakage} (Fuga de Datos) en el preprocesamiento manual.
        \item Comprender el concepto y estructura de un \textbf{Pipeline} en \texttt{scikit-learn}.
        \item Automatizar la imputación, el escalado y la clasificación en una sola línea de ensamblaje.
    \end{itemize}
\end{frame}

% ---------------------------------------------------
\begin{frame}{El Problema: Procesamiento Manual}
    Hasta hoy, hemos construido modelos paso a paso:
    \vspace{0.3cm}
    \begin{enumerate}
        \item Tomar los datos crudos.
        \item Rellenar los valores nulos (Imputar).
        \item Estandarizar los datos (Escalar).
        \item Entrenar el modelo.
    \end{enumerate}
    \vspace{0.3cm}
    \textbf{¿Cuál es el riesgo?} Si aplicamos estas transformaciones de forma manual, es muy fácil olvidar un paso al evaluar nuevos datos, o peor aún, utilizar datos del futuro (Set de Prueba) para calcular promedios de imputación o escalado (\textbf{Data Leakage}).
\end{frame}

% ---------------------------------------------------
\begin{frame}{La Solución: Línea de Ensamblaje (Pipeline)}
    \begin{center}
        % Referencia a la gráfica generada por generar_graficas.py
        \includegraphics[width=\textwidth,height=0.7\textheight,keepaspectratio]{pipeline_diagrama.png}
    \end{center}
\end{frame}

% ---------------------------------------------------
\begin{frame}{Anatomía de un Pipeline (Paso a Paso)}
    Un Pipeline es un objeto que encapsula múltiples pasos secuenciales.
    \vspace{0.3cm}
    \begin{itemize}
        \item \textbf{Estaciones (Transformers):} Cada paso intermedio debe tener las funciones \texttt{fit()} y \texttt{transform()} para modificar los datos (ej. Imputador y Escalador).
        \item \textbf{Cierre (Estimator):} El último paso de la cadena debe ser obligatoriamente un modelo predictivo, el cual posee la función \texttt{predict()} (ej. Árbol de Decisión).
    \end{itemize}
\end{frame}

% ---------------------------------------------------
\begin{frame}[fragile]{Implementación en Python: Importación de Herramientas}
    Primero, importamos todas las estaciones de nuestra fábrica:
    \vspace{0.3cm}
    \begin{verbatim}
from sklearn.pipeline import Pipeline
from sklearn.impute import SimpleImputer
from sklearn.preprocessing import StandardScaler
from sklearn.tree import DecisionTreeClassifier
    \end{verbatim}
    \vspace{0.3cm}
    Cada una de estas clases representa una máquina independiente que debemos ensamblar.
\end{frame}

% ---------------------------------------------------
\begin{frame}[fragile]{Implementación en Python: Ensamblaje}
    Se declara el Pipeline utilizando una lista de tuplas. Cada tupla contiene un nombre identificador y la herramienta instanciada.
    \vspace{0.3cm}
    \begin{verbatim}
mi_fabrica = Pipeline(steps=[
    ('limpiador', SimpleImputer(strategy='mean')),
    ('escalador', StandardScaler()),
    ('modelo', DecisionTreeClassifier(max_depth=4))
])
    \end{verbatim}
    \vspace{0.3cm}
    \textbf{Punto Crítico:} El orden es estrictamente secuencial. Los datos crudos entrarán por 'limpiador' y saldrán por 'modelo'.
\end{frame}

% ---------------------------------------------------
\begin{frame}[fragile]{Implementación en Python: Encendido Total}
    Una vez ensamblado el Pipeline, todo el flujo de trabajo se simplifica drásticamente:
    \vspace{0.3cm}
    \begin{verbatim}
# Todo el preprocesamiento y entrenamiento ocurre aquí
mi_fabrica.fit(X_train, y_train)

# El preprocesamiento se aplica automáticamente a los 
# datos de evaluación antes de predecir
predicciones = mi_fabrica.predict(X_test)
    \end{verbatim}
    \vspace{0.3cm}
    Esta estructura garantiza que el modelo puede ser puesto en producción y recibir datos crudos de forma segura.
\end{frame}

% ---------------------------------------------------
\begin{frame}[fragile]
    \begin{center}
        \Large \textbf{Fin de la Presentación Teórica}\\
        \vspace{0.5cm}
        \normalsize Transición al Mini Proyecto:\\
        \vspace{0.2cm}
        \verb|06_Pipelines_Mini_Proyecto.ipynb|\\
        \vspace{0.2cm}
        \verb|06_Laboratorio_Evaluacion.ipynb|
    \end{center}
\end{frame}

\end{document}
